%; whizzy paragraph -pdf xpdf -latex ./whizzypdfptex.sh
%; whizzy-paragraph "^\\\\begin{frame}\\|\\\\emtext"
% latex beamer presentation.
% platex, latex-beamer でコンパイルすることを想定。

%     Tokyo Debian Meeting resources
%     Copyright (C) 2012 Junichi Uekawa

%     This program is free software; you can redistribute it and/or modify
%     it under the terms of the GNU General Public License as published by
%     the Free Software Foundation; either version 2 of the License, or
%     (at your option) any later version.

%     This program is distributed in the hope that it will be useful,
%     but WITHOUT ANY WARRANTY; without even the implied warranty of
%     MERCHANTABILITY or FITNESS FOR A PARTICULAR PURPOSE.  See the
%     GNU General Public License for more details.

%     You should have received a copy of the GNU General Public License
%     along with this program; if not, write to the Free Software
%     Foundation, Inc., 51 Franklin St, Fifth Floor, Boston, MA  02110-1301 USA

\documentclass[cjk,dvipdfmx,12pt]{beamer}
\usetheme{Tokyo}
\usepackage{monthlypresentation}

%  preview (shell-command (concat "evince " (replace-regexp-in-string "tex$" "pdf"(buffer-file-name)) "&"))
%  presentation (shell-command (concat "xpdf -fullscreen " (replace-regexp-in-string "tex$" "pdf"(buffer-file-name)) "&"))
%  presentation (shell-command (concat "evince " (replace-regexp-in-string "tex$" "pdf"(buffer-file-name)) "&"))

%http://www.naney.org/diki/dk/hyperref.html
%日本語EUC系環境の時
\AtBeginDvi{\special{pdf:tounicode EUC-UCS2}}
%シフトJIS系環境の時
%\AtBeginDvi{\special{pdf:tounicode 90ms-RKSJ-UCS2}}

\newenvironment{commandlinesmall}%
{\VerbatimEnvironment
  \begin{Sbox}\begin{minipage}{1.0\hsize}\begin{fontsize}{8}{8} \begin{BVerbatim}}%
{\end{BVerbatim}\end{fontsize}\end{minipage}\end{Sbox}
  \setlength{\fboxsep}{8pt}
% start on a new paragraph

\vspace{6pt}% skip before
\fcolorbox{dancerdarkblue}{dancerlightblue}{\TheSbox}

\vspace{6pt}% skip after
}
%end of commandlinesmall


\title{東京エリアDebian勉強会\\　\\Buffalo Linkstation向け\\Debian Installer}
\subtitle{第139回 2016年5月度}
\author{Roger Shimizu}
\date{2016年5月21日}
\logo{\includegraphics[width=8cm]{image200607/openlogo-light.eps}}

\begin{document}

\frame{\titlepage{}}

\section{Agenda}
\begin{frame}{Agenda}
\begin{enumerate}
\item 自己紹介
\item 始めに
\item Buffalo Linkstation と Debian Installer
\item Linkstation へ Debian のインストール仕方
\item 終わりに
\item Demo
\end{enumerate}
\end{frame}

\section{自己紹介}
\emtext{自己紹介}

\begin{frame}{自己紹介}
\begin{itemize}
\item 10年以上(静かな) Debian user.
\item 2年前にLinkstation LS-WXL の Device-tree 対応を初め、Linux Kernel へ ARM Device の Device-Tree や Debian Project など、様々の貢献を行いました。
\item Debian の wide-dhcpv6, adjtimex パッケージをメインテインしております。
\item Debian Project の Kernel Team と Installer Team を参加しており、組み込み系 Device がより簡単に使えるように頑張ります。
\end{itemize}
\end{frame}

\section{始めに}
\emtext{始めに}

\begin{frame}{始めに}
VGA/HDMI port や Keyboard など付いてない機器、例えばサーバ向けてPC、それから ARM 開発ボードなどの機器に、どうやって Debian を入れておくでしょうか？
RAID 付きの Buffalo Linkstation NAS を例として、紹介致します。
\end{frame}

\section{Buffalo Linkstation と Debian Installer}
\emtext{Buffalo Linkstation と Debian Installer}

\begin{frame}{Buffalo Linkstation NAS の歴史}
\begin{itemize}
\item 第０世代: Linkstation / Kuro-Box
\item 第１世代: Linkstation HG / Kuro-Box HG
	\begin{itemize}
	\item PowerPC architecture
	\item IDE/PATA のみ (SATA がなし)
	\item 現在あっても使いづらいと思います
	\end{itemize}
\item 第２世代\footnote{MIPSのModelも出たけれど、Model数が少ないし、出てすぐデスコンになってしまったので、こちら省略させて頂きます。}: 
Kuro-Box Pro / Linkstation Live\\/ Linkstation LS-GL/LS-WTGL/LS-WSGL/LS-QL
	\begin{itemize}
	\item ARM architecture
		\begin{itemize}
		\item Debian Etchまで: arm OABI (Old ABI)
		\item Debian Lennyから: armel EABI (new Embedded ABI)
		\end{itemize}
	\item Marvell orion5x 5182 chipset; SATA interface
	\end{itemize}
\end{itemize}
\end{frame}

\begin{frame}{Buffalo Linkstation NAS 歴史 (続き)}
\begin{itemize}
\item 第３世代: Linkstation LS-XHL/LS-CHL/LS-WXL\\
\& Linkstation LS-VL/LS-WVL/LS-QVL, etc
	\begin{itemize}
	\item Marvell kirkwood 6281 / 6282 chipset
	\item armel architecture なので、第２世代と rootfs の互換性があり、また Debian Kernel 4.4 から同じ「-marvell」との flavour で共通に対応されます。
\footnote{HDDを異なる型番の機器に入れる前に、flash-kernelで適切なDTBをuImage.buffaloに入れ置かないと行けません。}
	\end{itemize}
\item 第４世代: LS-210/LS-220/LS-410/LS-420, etc
	\begin{itemize}
	\item Marvell armada-370 chipset
	\item armhf architecture (hard-float)
	\item 残念ながら、Linux Kernel の方がまだ対応されてないようです。
	\end{itemize}
\end{itemize}
\end{frame}

\begin{frame}{Debian Installer の紹介}
\begin{itemize}
\item Debian Installer (略は D-I となります)では、様々機器に Debian をインストールしてくれるツールとなります。
	\begin{itemize}
	\item 非 Linux な OS でも対応される、例え kFreeBSD や GNU/Hurd など
	\item メディアは CD/DVD に限られず、PXE netboot や u-boot など様々柔軟なインストールメディアを提供されております。
	\item モニターなど表示デバイスが付いてない機器でも対応される──network-console イメージ\\
	　⇒ ちょうど Buffalo Linkstation など NAS 機器に向け
	\end{itemize}
\end{itemize}
\end{frame}

\begin{frame}{Debian Installer の紹介 (続き)}
\begin{itemize}
\item 現在 D-I に対応されている Linkstation リスト\footnote{d-i daily image が含まれる}:
	\begin{itemize}
	\item Kuro-Box Pro / Linkstation Pro/Live
	\item Linkstation LS-GL / LS-WSGL / LS-WTGL
	\item Linkstation LS-XHL / LS-CHLv2 / LS-WXL / LS-WSXL / LS-VL / LS-WVL / LS-QVL\\
	　⇒ armel の第２世代と第３世代はほとんど対応されてます
	\end{itemize}
\end{itemize}
\end{frame}


\section{Linkstationへ Debian の インストール仕方}
\emtext{Linkstationへ Debian の インストール仕方}

\begin{frame}{Linkstation に Debian Installer を起動させる}
Buffalo Linkstation では、u-boot という boot loader が、1番目の partition (/dev/sda1 又は /dev/md0) に保存される kernel と initrd を読み込んで起動を行います。
それから、D-I のイメージを1番目の partition に置いとけば、Debian Installer が起動されます。
手順は下記の通りとなります。
\begin{itemize}
\item 1番目の partition をフォーマットする\footnote{1番目の partition は ext2/ext3 に限られます。}(既に存在するなら省略可)
\item D-I kernel/initrd イメージを、1番目の partition にコピーする。D-I イメージは下記URLでダウンロード出来ます。
	\begin{itemize}
	\item https://d-i.debian.org/daily-images/armel/daily\\/orion5x/network-console/buffalo
	\item https://d-i.debian.org/daily-images/armel/daily\\/kirkwood/network-console/buffalo
	\end{itemize}
\end{itemize}
\end{frame}

\begin{frame}{Linkstation に Debian Installer を起動させる(続)}
\begin{itemize}
\item D-I kernel/initrd image を入れた HDD を Linkstation に取り付けます。
\item Linkstation を起動し、DHCP で IP Address を割り当てられるまで暫く待って、それから割り当てられた IP Address に SSH で接続します。
\item 画面に沿って、通常な Debian Install 行うことが出来ます。
\end{itemize}
\end{frame}

\begin{frame}[containsverbatim]{詳細な手順: HDD Partition の作成}
LinkstationのHDDをLinkstationから外し、SATA-USBアダプターなど方法経由でPCに接続します。例えば、その HDD が /dev/sdc として説明します。
LS-WXL/WSXL/WVL/QVL など RAID 機器の場合は、すべての HDD を同じ partition にしても良いです。
\begin{commandline}
$ sudo parted /dev/sdc
(parted) mklabel gpt
(parted) mkpart boot 2048s 1024MiB
(parted) mkpart root 1024MiB 6144MiB
(parted) mkpart swap 6144MiB 6400MiB
(parted) mkpart data 6400MiB -1
# 下記のコマンドは RAID 構成の機種だけ必要となります。
(parted) set 1 raid on
(parted) set 2 raid on
(parted) set 3 raid on
(parted) set 4 raid on
\end{commandline}
\end{frame}

\begin{frame}[containsverbatim]{詳細な手順: HDD Partition の確認}
\begin{itemize}
\item Partition を作ったら、確認するとこうなります。
\end{itemize}
\begin{commandline}
(parted) print
Model: SAMSUNG HM250HI (scsi)
Disk /dev/sdc: 250GB
Sector size (logical/physical): 512B/512B
Partition Table: gpt
Disk Flags:
Num Start   End     Size    File sys  Name  Flags
 1  1049kB  1074MB  1073MB            boot  raid
 2  1074MB  6442MB  5369MB            root  raid
 3  6442MB  6711MB  268MB             swap  raid
 4  6711MB  250GB   243GB             data  raid

# 最後にpartedを終了させます。
(parted) quit
\end{commandline}
\begin{itemize}
\item RAID 構成の場合は他の HDD も同じようにセットしてください。
\end{itemize}
\end{frame}

\begin{frame}[containsverbatim]{詳細な手順: boot image のセット}
\begin{itemize}
\item boot image kernel/initrd を第１番目 partition にコピーします。(例え、LS-WXL のイメージ)
\item RAID の複数HDDの場合は、念のためすべてHDDをコピーしましょう。\footnote{この時点に「mdadm --create」で RAID の構成をしなくても良いです。Installの際に RAID の作成を行います。}
\end{itemize}
\begin{commandline}
$ sudo mkfs.ext3 /dev/sdc1
$ sudo mount /dev/sdc1 /mnt
$ wget https://d-i.debian.org/daily-images/armel\
  /daily/kirkwood/network-console/buffalo/ls-wxl\
  /uImage.buffalo
$ wget https://d-i.debian.org/daily-images/armel\
  /daily/kirkwood/network-console/buffalo/ls-wxl\
  /initrd.buffalo
$ sudo cp *.buffalo /mnt
$ sudo umount /mnt
\end{commandline}
\end{frame}

\begin{frame}[containsverbatim]{詳細な手順: SSH で接続}
\begin{itemize}
\item HDDs を Linkstation に戻して、起動させます。
\item 暫く待つと、Android/iOS アプリ「Fing」のような IP/port scanner で Linkstation に割り当てた IP を見つけます。また、DHCP サーバ側のログでも
\item IP を分かると、SSH コマンドを叩くと debian installer 画面が出て来る：
\end{itemize}
\begin{commandline}
$ ssh installer@<IP address of Linkstation>
\end{commandline}
\begin{itemize}
\item D-I のデフォルトパスワードは「install」となります。
\item command line で操作や log 確認などのため、もう一本 SSH を接続しても良いです。
\end{itemize}
\end{frame}

\begin{frame}[containsverbatim]{RAID 構成向け Install 時の注意事項 (その０)}
\begin{itemize}
\item LS-GL/CHL/XHL/VL などRAID構成ではない機種だとスキップください。
\item もし RAID の設定が見つからない場合は、D-I Partman の画面から一旦「バック」し、「Download installer components」に partman-md や sata-modules などモジュールを選択してから、出来るようになります。
\end{itemize}
\end{frame}

\begin{frame}[containsverbatim]{RAID 構成向け Install 時の注意事項 (その１)}
\begin{itemize}
\item もし D-I で RAID を新規に作成する場合、一番目の /dev/md0 は metadata=0 (version 0.90) に設定しないと再起動しなくなります。原因は u-boot が第１番目 partition の kernel/initrd を読み込まなくなるため。現在 partman-md に設定出来なくて\footnote{\url{https://bugs.debian.org/815569}}、command line にしましょう：
\begin{commandline}
# mdadm --create /dev/md0 --level=1 --raid-devices=2 \
  --metadata=0 /dev/sda1 /dev/sdb1
\end{commandline}
\item または (他の HDD が後にします)
\begin{commandline}
# mdadm --create /dev/md0 --level=1 --raid-devices=2 \
  --metadata=0 /dev/sda1 missing
\end{commandline}
\item /dev/md0 以外の RAID は partman-md で設定しても良いです。
\end{itemize}
\end{frame}

\begin{frame}[containsverbatim]{RAID 構成向け Install 時の注意事項 (その２)}
\begin{itemize}
\item RAID を新規に作成する (create) と同期化の作業がすぐに始めされ、とても重くて、進行中 Debian Install の作業に影響ならないように /dev/md0 以外の同期化の速度を制限かけた方が良いです：
\end{itemize}
\begin{commandline}
# echo 100 > /sys/block/md{1,2,3}/md/sync_speed_max
\end{commandline}
\begin{itemize}
\item インストールが終わって再起動したら、RAID の同期化の作業は自動的に再開されるので、その後の作業は特に必要ありません。
\end{itemize}
\end{frame}

\begin{frame}[containsverbatim]{Debian Install 後の設定}
\begin{itemize}
\item u-boot 環境変数を確認・変更のため、コマンド fw\_printenv / fw\_setenv の設定
\footnote{u-boot 環境変数を変更すると起動しない場合があり、修復手段があまりなくて、気をつけましょう。}
	\begin{itemize}
	\item ほとんどの機種は下記の設定が済みです：
	\end{itemize}
\begin{commandline}
$ sudo echo /dev/mtd2 0x00000 0x10000 0x10000 \
	> /etc/fw_env.config
\end{commandline}
	\begin{itemize}
	\item Kuro-Box Pro は mtd/flash の数が異なるので、下記の設定にしてください。
	\end{itemize}
\begin{commandline}
$ sudo echo /dev/mtd5 0x00000 0x10000 0x10000 \
	> /etc/fw_env.config
\end{commandline}
\end{itemize}
\end{frame}

\begin{frame}[containsverbatim]{Debian インストール後の設定 (続き)}
\begin{itemize}
\item 他の機器で Linkstation の起動ログを確認するため(あるいは、再起動不能の際にデバッグ手段として)、netconsole の設定
	\begin{itemize}
	\item Linkstation 側の設定：
	\end{itemize}
\begin{commandline}
$ sudo cat << EOT >> /etc/initramfs-tools/modules
marvell
mv643xx_eth
netconsole netconsole=@192.168.11.5/,6666@192.168.11.1/
mvmdio
EOT

$ sudo update-initramfs -u
\end{commandline}
	\begin{itemize}
	\item 他の機器でログ収集するコマンド：
	\end{itemize}
\begin{commandline}
$ sudo ip a add 192.168.11.1/24 dev eth0
$ nc -l -u -p 6666 | tee ~/netconsole.log
\end{commandline}
\end{itemize}
\end{frame}


\section{終わりに}
\emtext{終わりに}

\begin{frame}{終わりに}
\begin{itemize}
\item Kernel の Device-Tree を対応したり、Debian-Installer にパッチを投げたり、Linkstation に Debian Install はやっと出来るようになりました。
\item 今後は Debian Installer を引き続き改善・進化を行って行きたいと思います。
	\begin{itemize}
	\item Debian Installer に GNU/screen 又は tmux を対応する
	\footnote{\url{https://lists.debian.org/debian-boot/2016/02/msg00547.html}}\footnote{\url{https://bugs.debian.org/819397}}\footnote{\url{https://bugs.debian.org/819988}}
		\begin{itemize}
		\item SSH 接続が切れても、installer がバックグラウンドに回せて、SSH で再接続すると元の状態に再開できます。
		\item シェルやログなどより便利にアクセス出来ます。
		\end{itemize}
	\item partman-md へ RAID の metadata 指定できるように、対応する\footnote{\url{https://bugs.debian.org/815569}}。
	\item 第４世代、armhf/armada-370 の Linkstation を対応する
	\end{itemize}
\end{itemize}
\end{frame}


\section{Demo}
\emtext{Demo}

\begin{frame}{Buffalo Linkstation 向け Debian Install Demo}
\begin{itemize}
\item 機材: 2-Bay 2.5’ HDD model Linkstation LS-WSXL
\end{itemize}
\end{frame}

\end{document}

;;; Local Variables: ***
;;; outline-regexp: "\\([ 	]*\\\\\\(documentstyle\\|documentclass\\|emtext\\|section\\|begin{frame}\\)\\*?[ 	]*[[{]\\|[]+\\)" ***
;;; End: ***
